\section{Methods}
Here, we outline the technical foundations of our `Invariant Image Reparameterisation' approach. Although we combine ideas in a new way, we build in an important way on the frameworks reviewed in~\cite{Cole2010}, \cite{Cole2020}, \cite{Simpson2023}, and, indirectly,~\cite{Maclaren2019, Maclaren2021}.

\subsection{Model representation: auxiliary mapping}
With mechanistic models, it is useful to distinguish between underlying theoretical model parameters and parameters more directly associated with observable data or data distributions. The main motivation for this distinction is that there may be a more direct connection between data and data distribution parameters than between data and underlying mechanistic model parameters, even though the latter are the focus of identifiability analysis. We use $\theta \in \Theta \subseteq \mathbb{R}^p$ to denote the model parameters, and $y \in \mathcal{Y} \subseteq \mathbb{R}^n$ to denote observed (or predicted) data. The data may be scalar or vector-valued, and may represent spatial or temporal observations. We then assume a statistical model for the data, $p(y;\phi)$, where $p$ is a probability density function (or mass function) parameterised by $\phi$. We call $\phi \in \Phi \subseteq \mathbb{R}^q$ the data distribution parameters, which we assume are sufficient to characterise the empirical data distribution independently of any underlying mechanistic model. This may characterise a statistical observation model on top of a deterministic mechanistic model, but may also represent a purely statistical model capturing key features of the data assumed to be generated by some (potentially unknown) underlying and stochastic mechanism.

This general approach has been used in the context of econometrics, e.g., indirect inference and the method of moments~\cite{Gourieroux1993, Heggland2004}, and has also been used in the context of mechanistic models in biology and engineering~\cite{Drovandi2015, Browning2020, Simpson2021, Simpson2023}. A simple example occurs in the method of moments, where the data distribution parameters are the mean and variance of the data distribution, while the underlying model parameters may be more complex. There are also close connections to the `exhaustive summary' approach reviewed by~\cite{Cole2010, Cole2020}, where our data distribution parameters correspond to the exhaustive summaries.

Thus, the role of the model becomes to connect the parameters to data via data distribution parameters. Importantly, this is assumed to result in a well-defined deterministic mapping between these two parameter types, and so can be directly analysed as a function, though it may require numerical computation. For example, this mapping may require the numerical solution of differential equations or other complex computations. We call the mapping so defined the `auxiliary mapping' $\phi(\theta)$~\cite{Simpson2023}, which may be implicit or explicitly defined by the model. This gives a collection of relationships beginning from the auxiliary mapping:
\begin{equation}
\theta \mapsto \phi(\theta),
\end{equation}
mapping from model parameters to data distribution parameters. These data distribution parameters characterise the density function
\begin{equation}
p(y; \phi).
\end{equation}
Combined, we obtain the probability model parameterised in terms of model parameters:
\begin{equation}
p(y;\theta) = p(y;\phi(\theta)).
\end{equation}
For deterministic spatial/temporal mechanistic models considered here, we use numerical solutions on a fine grid as the data distribution parameters, or `exhaustive' summary of the mechanistic parameters in terms of predictions, evaluating derivatives via automatic differentiation in Julia~\cite{Revels2016}. Similarly, we might use e.g., the full trajectories of the moments defined by a moment dynamics approximation of a stochastic model~\cite{Browning2020}. Thus, while our approach is essentially identical to the `exhaustive summary' approach~\cite{Cole2010, Cole2020}, we generally avoid the need for explicit exhaustive summary methods such as Taylor series or generating series, though these can also be used if desired. In this context, we emphasise that automatic differentiation at the source code level is distinct from symbolic differentiation and only evaluates the numerical value of the derivative at a fixed input value~\cite{Griewank2008}. Numerical solutions are typically more readily available and applicable than symbolic methods for complex models.

\subsection{Model reparameterisation: basic ideas}
A key idea in identifiability analysis is that non-identifiable parameters can often be combined into identifiable parameter combinations or, more generally, that some functions of the parameters can be made identifiable~\cite{Maclaren2019,Cole2020}. This motivates the idea of reparameterising the model in terms of these identifiable combinations, which can be considered a type of (exact) model reduction. While this is often done symbolically, e.g.~\cite{Catchpole1997,Catchpole1998,Cole2010,Cole2020}, we show how this can be done numerically under certain conditions. We also show how this can be used to motivate reparameterisations corresponding to approximate model reduction in the case of practical non-identifiability. We also make an important distinction between what we call image reparameterisations, which contain all identifiable combinations but may also contain non-identifiable combinations, and minimal image reparameterisations, which contain only identifiable combinations. In particular an image reparameterisation may represent a reduced-dimension parameterisation, but this may still be redundant and hence not be identifiable as a whole. This is essentially the same ideas as that of exhaustive summaries, though here we use the term image reparameterisation in the context of our particular approach. The main use of an image reparameterisation is as an intermediate step towards a minimal image reparameterisation.

\subsection{Reparameterisations of the auxiliary mapping: existence of image and minimal image reparameterisations}
The core idea of our approach is that the auxiliary mapping can be decomposed into two parts: a reparameterisation that extracts the identifiable parameter combinations, followed by a one-to-one mapping of these to the data distribution parameters. We can write this as:
\begin{equation}
\phi(\theta) = \tilde{\phi}(\psi(\theta)),
\label{eq:epi-mono}
\end{equation}
where $\psi(\theta)$ represents identifiable parameter combinations defining the minimal image reparameterisation, and $\tilde{\phi}$ provides a one-to-one mapping from identifiable parameter combinations to data distribution parameters. The existence of such a minimal image reparameterisation is guaranteed by the epi-mono factorisation property of the general category of sets and mappings~\cite{Lawvere2003} (see supplementary material for details).

Unfortunately, this is a fairly abstract result and does not directly provide a method for determining the intermediate image or associated reparameterisation. To make this concept practical we assume our mappings are sufficiently smooth and use the language of differential geometry. In particular, we assume that the reparameterisation $\psi(\theta)$ is a smooth submersion (its derivative is surjective/onto) and the reduced mapping $\tilde{\phi}$ is a smooth immersion (its derivative is injective/one-to-one)~\cite{Lee2013}. This structure is a natural assumption for well-behaved models and, as shown in the next section, leads to a set of solvable symbolic conditions. Figure~\ref{fig:epi-mono} illustrates this decomposition.

\begin{figure}[htbp]
    \centering
    \includegraphics[width=0.6\linewidth]{epi-mono-factor.pdf}
    \caption{The basic decomposition of the auxiliary mapping into an identifiable reparameterisation followed by a reduced model mapping. Here $\theta$ represent the original model parameters, $\psi(\theta)$ represents identifiable parameter combinations (the minimal image reparameterisation), and $\tilde{\phi}$ provides a one-to-one mapping from identifiable parameter combinations to data distribution parameters.}
    \label{fig:epi-mono}
\end{figure}

Although quite general and solvable in principle using symbolic methods, solving these conditions symbolically is often intractable for complex models. The main contribution of our work is to convert this global symbolic problem into efficient local numerical computations that have global, or at least approximately global, validity. To do this, we introduce structural assumptions about the form of the reparameterisation. In particular, we consider cases where the model can be expressed in terms of monomial parameter combinations and, more generally, as functions of linear combinations of monomial parameter combinations. This is a natural assumption for many mechanistic models, e.g., via dimensional analysis~\cite{Buckingham1914, Barenblatt2003}, and also arises in asymptotic approximations and model reduction~\cite{Maiwald2016}. We will see that this assumption allows us to convert the symbolic reparameterisation condition to an equivalent numerical condition that can be evaluated at a single reference point, and hence avoid symbolic computation.

\subsection{Symbolic reparameterisation condition}
For smooth mappings with constant rank, we can strengthen the assumptions about the minimal image reparameterisation (epi-mono factorisation) by requiring $\psi(\theta)$ to be a smooth submersion and $\tilde{\phi}(\psi)$ to be a smooth immersion~\cite{Lee2013}. The chain rule then gives:
\begin{equation}
D_\theta\phi(\theta) = D_\psi\tilde{\phi}(\psi(\theta))D_\theta\psi(\theta).
\end{equation}
Since $D_\psi\tilde{\phi}(\psi(\theta))$ is one-to-one, the rank of $D_\psi\tilde{\phi}(\psi(\theta))$ and $D_\theta\phi(\theta)$ are the same; furthermore, these derivatives share the same null space/kernel and complementary row spaces (see supplementary material):
\begin{equation}
\begin{aligned}
\nullspace D_\theta\psi(\theta) &= \nullspace D_\theta\phi(\theta)\\
\rowspace D_\theta\psi(\theta) &= \rowspace D_\theta\phi(\theta).
\end{aligned}
\end{equation}
Given null space vectors $\alpha(\theta)$ obtained from $D_\theta\phi(\theta)$, we can then solve the differential condition:
\begin{equation}
D_\theta\psi(\theta)\alpha(\theta) = 0,
\end{equation}
for $\psi(\theta)$. This is equivalent to the conditions given by e.g.~\cite{Catchpole1997,Catchpole1998}, summarised by ~\cite{Cole2010, Cole2020}. We can also directly equate the rows of $D_\theta\psi(\theta)$ to the vectors spanning the row space of $D_\theta\phi(\theta)$, $\rowspace D_\theta\phi(\theta)$, which gives a system of differential equations for $\psi$. Abstractly, the image of the parameter space $\Theta$ under the resulting $\psi$ gives the intermediate `image' component of the epi-mono factorisation. A concrete example of solving the above symbolic condition is provided in the Results section.

The above provides conditions to solve for a minimal image reparameterisation. In some cases we may seek or already have a proposed image reparameterisation that is not minimal. An image reparameterisation that is not necessarily minimal just means that we can write the auxiliary mapping as:
\begin{equation}
\phi(\theta) = \tilde{\phi}(\psi(\theta)),
\label{eq:suff-reparam}
\end{equation}
without the guarantee that $\tilde{\phi}$ is one-to-one. This provides a reduced-dimension parameterisation that by definition doesn't lose information in the sense that the model outputs are the same, but this parameterisation may still be redundant. In such a case we cannot equate the null spaces directly, however we can show that the following weaker condition holds (see supplementary material):
\begin{equation}
    \begin{aligned}
\nullspace D_\theta\psi(\theta) &\subseteq \nullspace D_\theta\phi(\theta)\\
\rowspace D_\theta\psi(\theta) &\supseteq \rowspace D_\theta\phi(\theta).
\end{aligned}
\end{equation}
This concept has two important applications in the symbolic context. First, given a proposed simplification $\psi(\theta)$ from theory or other sources, this condition serves as a test to verify that it is a valid, lossless reduction. Second, one can derive an image reparameterisation satisfying the above conditions by enforcing the orthogonality condition, $D_\theta\psi(\theta)\alpha(\theta)=0$, for a set of vectors $\{\alpha_i\}$ that span only a proper subspace of the full model's null space. We discuss the choice of an appropriate subspace, and the utility of this condition in the context of numerical approaches, below.

The downside of these symbolic approaches is that they must be solved symbolically, i.e., enforced for all $\theta$ (in both the image and minimal image cases), which requires symbolic computation and is typically only possible for simple models. We again emphasise that this symbolic approach is essentially equivalent to the approaches reviewed by~\cite{Cole2010, Cole2020}, building on e.g. \cite{Catchpole1997,Catchpole1998}, though we use slightly different terminology and notation. We next show how this can be converted to a numerical approach under certain conditions.

\subsection{Reparameterisation: structural assumptions and parameter space linearisation}
Our goal is to determine identifiable parameter combinations without symbolic computation. The primary obstacle to a numerical identifiability analysis is that local derivative information is, by definition, parameter-dependent. Thus we want to find equivalent constraints on our reparameterisation that are independent of $\theta$, so that we can evaluate them at a single convenient reference point. We achieve this by restricting the class of reparameterisations we consider to those having a particular structural form.

\subsubsection{Single-stage transformation}
In particular, we assume a reparameterisation with the structural form:
\begin{equation}
\psi(\theta) = f^{-1} \circ A \circ f(\theta),
\end{equation}
where $A$ is a constant (parameter-independent), `wide' (number of columns $\geq$ number of rows) but otherwise arbitrary matrix (linear transformation) and $f$ is an element-wise nonlinear (in general), invertible, and strictly monotonic (in each entry of $\theta$) map. The matrix $A$ will not in general be invertible. The constant (invariant) nature of $A$ is the key property that allows a local, numerical computation to have global, or at least approximately global, validity, and we aim to equate this to parameter-independent (invariant) properties of the model Jacobian in appropriate coordinates.

This structure follows a conjugacy-type pattern: we transform parameters to a new coordinate system via $f$, perform a linear transformation $A$ in that space (where Algorithm~\ref{alg:reparam} applies standard linear analysis via Jacobian and Hessian computations), then transform back via $f^{-1}$ to express results in the original parameter space. This is analogous to conjugacy as used in dynamical systems to establish topological equivalence between a nonlinear system and its linearisation~\cite{Wiggins2003}. Though a different application, we similarly use this structure to `exactly linearise' some aspects of our system without losing information.

A case of particular interest is the form:
\begin{equation}
\psi(\theta) = \exp \circ A \circ \log (\theta) = \exp(A \log \theta),
\end{equation}
which generates a vector of monomial combinations of the entries of an input vector, i.e., each entry of $\psi(\theta)$ is a monomial of the form
\begin{equation}
\psi_j(\theta) = \theta_1^{A_{j1}}\theta_2^{A_{j2}}\cdot\dots\theta_p^{A_{jp}}
\end{equation}
where $A_{ij}$ are the entries of $A$. These powers are integers for strict monomials, though we allow non-integer values in general. Here we assume that the sign of each entry of $\theta$ is either known or can be estimated (via the maximum likelihood estimate, for example). Hence we can apply the log function to the absolute value of each entry and then reintroduce the sign in the auxiliary mapping if needed. Log parameter transformations are also commonly used in the sloppiness literature~\cite{Brown2004,Monsalve2022,Vollert2023} and this also partly inspired our present approach, though our approach explicitly treats the log transformation as part of general structural reparameterisation transformations.

\subsubsection{Sequential multi-stage transformations}
We extend this analysis to sequences of such transformations applied iteratively. Each stage follows the same conjugacy-type pattern (transform to new coordinates, linear analysis, transform back), but with different choices of the transformation function $f$. The choice of $f$ at each stage depends on the anticipated structure:
\begin{itemize}
\item $f = \log$: Maps products to sums, exposing multiplicative/log-monomial combinations
\item $f = \text{identity}$: No transformation, exposing additive/linear combinations directly
\end{itemize}

For example, a two-stage transformation results in:
\begin{equation}
\psi = A_2 \circ \exp \circ A_1 \circ \log,
\end{equation}
for wide but otherwise arbitrary matrices $A_1$ and $A_2$. This form generates a broader class of outputs consisting of linear combinations of monomials. We approach this by applying exactly the same algorithm (Algorithm~\ref{alg:reparam}) recursively: the first stage uses $f=\log$ to identify monomial combinations, creating new coordinates; the second stage uses $f=\text{identity}$ on these coordinates to identify linear combinations among them. This sequential approach can be viewed as a series of image reparameterisations, with the final stage potentially yielding a minimal image reparameterisation in terms of the parameter combinations generated by the earlier stages. The equality or inclusion conditions from the previous section apply at each stage separately.

While two stages (log followed by identity) suffice for the examples considered here, the framework naturally extends to additional stages with other choices of $f$ for more complex compositional structures. This compositional approach enables discovery of mixed structures like linear combinations of monomials (e.g., $n_1p_1 + n_2p_2$) without symbolic computation. This multi-stage form is also analogous to a multiple-layer neural network with fixed nonlinearities~\cite{Goodfellow2016}, though again our motivation and assumptions are different.

\subsubsection{Reparameterisation condition: general form arising from the basic structural assumption}
We first consider the general single-stage reparameterisation of the form $\psi(\theta) = f^{-1}(A f(\theta))$. We define $\theta^* = f(\theta)$ to be the component-wise $f$ applied to the parameter vector, and define $\phi_*(\theta^*) = \phi(\theta)$ and $\psi_*(\theta^*) = \psi(\theta)$ to be the corresponding mappings in these `$f$-coordinates'. For example, $f$ may be the component-wise log function as in the monomial case below. Models that can be written in terms of our reparameterisation $\psi(\theta)$ then take the form:
\begin{equation}
\phi(\theta) =\phi_*(\theta^*) = \tilde{\phi}(f^{-1}(A\theta^*)),
\label{eq:reparam}
\end{equation}
where $A$ is a matrix of coefficients and $\tilde{\phi}$ is the intermediate mapping in \eqref{eq:epi-mono} or \eqref{eq:suff-reparam}. The chain rule gives:
\begin{equation}
D_{\theta^*} \phi_*(\theta^*) = D_{\psi_*} \tilde{\phi}(\psi_*(\theta^*))\text{diag}(f^{-1'}(A\theta^*))A,
\label{eq:chainrule}
\end{equation}
where $f^{-1'}$ is the component-wise derivative of $f^{-1}$. Now, in the case that our reparameterisation is a minimal image reparameterisation, i.e. $f^{-1}$ and $f$ are both smooth immersions, we have (by the same argument as before) the key conditions providing a decomposition of the fundamental spaces of $A$:
\begin{equation}
    \begin{aligned}
\nullspace A &= \nullspace D_{\theta^*}\phi_*(\theta^*)\\
\rowspace A &= \rowspace D_{\theta^*}\phi_*(\theta^*).
\end{aligned}
\label{eq:key-condition}
\end{equation}
This follows since the first term on the right-hand side of the chain rule equation \eqref{eq:chainrule} is injective by assumption, and the second term is a diagonal matrix with non-zero entries (since $f^{-1'}$ is non-zero by assumption). Furthermore, as the left-hand side of \eqref{eq:key-condition} is parameter-independent, the right-hand side must be too. We can hence evaluate the right-hand side derivative at any convenient point to determine $A$, reducing what would be a PDE to an algebraic computation. If the model is assumed non-identifiable, the point used will generally be, e.g., any (non-unique) maximum likelihood estimate, and the Jacobian will be singular. This is expected and causes no issues in general.

When our reparameterisation is only an image but not a minimal image reparameterisation, and hence the left-most term in the right-hand side of the chain rule equation \eqref{eq:chainrule} is not injective, we still have the weaker conditions:
\begin{equation}
    \begin{aligned}
\nullspace A &\subseteq \nullspace D_{\theta^*}\phi_*(\theta^*)\\
\rowspace A &\supseteq \rowspace D_{\theta^*}\phi_*(\theta^*).
\end{aligned}
\end{equation}
We can turn this into a constructive approach for determining $A$ by enforcing the orthogonality condition on only a subspace of the full null space of $D_{\theta^*}\phi_*(\theta^*)$, as described above.

In both cases, we need to verify that we have an image or minimal image reparameterisation. In the former case, we need to decide which subspace of the full null space to enforce orthogonality for, in the latter we enforce them all but assume the associated null-space is constant. The key idea hence is to either assume or verify that the null space of $D_{\theta^*}\phi_*(\theta^*)$, or a null subspace of it, is constant. We discuss how to verify this in practice below. We first note that in the case of an image reparameterisation, this idea can clearly be applied sequentially: since the output of one reparameterisation is an image reparameterisation, we can again seek a reparameterisation starting with this parameterisation and using a new $f$ and $A$. This means we can build up a more complex reparameterisation such as linear combinations of monomials (see below).

While the above condition does not uniquely determine $A$, as any basis for the null space of the model Jacobian (and complementary row space) can be used, it establishes the existence of a matrix capturing non-identifiable and (potentially) identifiable parameter combinations. While the coefficients in $A$ are traditionally restricted to integers in dimensional analysis, here we allow them to be arbitrary real numbers and do not require the combinations to be dimensionless \textit{a priori}. This relaxation facilitates numerical computation and provides more flexibility in discovering practically useful parameter combinations. However, we also consider \textit{a posteriori} scaling and rounding operations to convert these real-valued coefficients to integers to aid interpretability while preserving key properties of the approximation.

The corresponding alternative route to the decomposition of the auxiliary mapping is illustrated in Figure~\ref{fig:epi-linear}.

\begin{figure}[htbp]
    \centering
    \includegraphics[width=0.6\linewidth]{epi-mono-linearised.pdf}
    \caption{Implementation of parameter space linearisation using an initial componentwise transformation $f(\theta)$, typically $\log$, followed by a linear transformation and componentwise inverse. The chain rule applied on the path from $\Theta^*$ to $\Phi$ allows us to equate the overall null space to that of the constant linear transformation $A$.}
    \label{fig:epi-linear}
\end{figure}


\subsubsection{Reparameterisation: invariant subspace conditions}
[Rest of section unchanged - the invariance test is the same regardless of which stage we're at]
...

\subsection{Overall algorithm}
The overall approach is summarised in the following algorithm.
\begin{algorithm}[H]
\caption{Numerical reparameterisation algorithm}
\label{alg:reparam}
\begin{algorithmic}[1]
\Input Auxiliary mapping $\phi(\theta)$, reference parameter value $\hat{\theta}$, component-wise transformations $f$ and $f^{-1}$ (e.g., $f=\log$, $f^{-1}=\exp$)
\Output Reparameterisation matrix $A$ for reparameterisation $\psi = f^{-1} \circ A \circ f$
\State Compute $\hat{\theta}^* = f(\hat{\theta})$
\State Compute the Jacobian matrix $D_{\theta^*}\phi_*(\hat{\theta}^*)$ using automatic differentiation or other method
\State Compute the SVD $D_{\theta^*}\phi_*(\hat{\theta}^*) = U\Sigma V^T$
\State Partition $V=[V_r, V_0]$ based on the non-zero and zero singular values
\State Form the stacked matrix $M_{\text{stack}} = [H_1 V_0; H_2 V_0; \dots; H_p V_0]$ appearing in the invariance condition \eqref{eq:invariance-condition}
\State Compute the SVD of $M_{\text{stack}}$ and partition the right singular vectors of $M_{\text{stack}}$ into $V_M=[V_{Mr}, V_{M0}]$ based on non-zero and zero singular values
\State Construct a basis for the complement, $N_{\perp} = [V_r, V_0V_{Mr}]$, of the invariant null space $N = V_0V_{M0}$
\State Set the reparameterisation matrix $A = N_{\perp}^T$
\State Optionally scale and round entries of $A$ to improve interpretability
\State Return $A$ for reparameterisation $\psi = f^{-1} \circ A \circ f$
\end{algorithmic}
\end{algorithm}
This algorithm provides a reparameterisation $\psi(\theta) = f^{-1} (A f(\theta))$ that is either a minimal image or just an image reparameterisation, depending on whether the full null space or only an invariant subspace is found (this can be seen from the dimensions of the various matrices involved). The reparameterised model can then be used for parameter estimation and uncertainty quantification. As noted above, the algorithm can be applied sequentially to build up more complex reparameterisations. In this case, the output of one application of the algorithm becomes the input to the next, with a new choice of $f$ and $A$. For example, to find linear combinations of monomials, the first application uses $f=\log$ and $f^{-1}=\exp$ to determine $A_1$, while the second application uses $f=\text{id}$ and $f^{-1}=\text{id}$ on the first-stage coordinates to determine $A_2$, giving the overall reparameterisation $\psi(\theta) = A_2 \exp(A_1 \log \theta)$.

[Rest of Methods section unchanged...]
